\documentclass[11pt,a4paper]{article}

% ====== Pacotes ======
\usepackage[utf8]{inputenc}
\usepackage[T1]{fontenc}
\usepackage[brazilian]{babel}
\usepackage{lmodern}
\usepackage[margin=2.4cm]{geometry}
\usepackage{amsmath,amssymb}
\usepackage{xcolor}
\usepackage[most]{tcolorbox}
\usepackage{enumitem}
\usepackage{hyperref}
\usepackage{fancyhdr}
\usepackage{titlesec}
\usepackage{microtype}
\usepackage{array}
\usepackage{booktabs}
\usepackage{listings}

% ====== Cores ======
\definecolor{deitelblue}{RGB}{31,73,125}
\definecolor{deitelorange}{RGB}{198,89,17}
\definecolor{boapratcolor}{RGB}{0,112,60}
\definecolor{errocolor}{RGB}{178,34,34}
\definecolor{engcolor}{RGB}{31,73,125}
\definecolor{desempcolor}{RGB}{198,89,17}
\definecolor{portabcolor}{RGB}{102,46,145}
\definecolor{codebg}{RGB}{248,248,248}
\definecolor{codekeyword}{RGB}{0,0,180}
\definecolor{codestring}{RGB}{163,21,21}
\definecolor{codecomment}{RGB}{0,128,0}
\definecolor{terminalbg}{RGB}{30,30,30}
\definecolor{terminalfg}{RGB}{220,220,220}

% ====== Listings ======
\lstdefinestyle{cppcode}{
  language=C++,
  basicstyle=\ttfamily\footnotesize,
  keywordstyle=\color{codekeyword}\bfseries,
  stringstyle=\color{codestring},
  commentstyle=\color{codecomment}\itshape,
  numbers=left,
  numberstyle=\tiny\color{gray},
  numbersep=8pt,
  backgroundcolor=\color{codebg},
  frame=single,
  framesep=4pt,
  rulecolor=\color{deitelblue!50},
  breaklines=true,
  showstringspaces=false,
  tabsize=4,
  morekeywords={string, size_t, cin, cout, endl, inline, template, typename,
                bool, true, false, nullptr, override, final},
  literate={ã}{{\~a}}1 {á}{{\'a}}1 {é}{{\'e}}1 {ê}{{\^e}}1 {í}{{\'i}}1
           {ó}{{\'o}}1 {ô}{{\^o}}1 {õ}{{\~o}}1 {ú}{{\'u}}1 {ç}{{\c{c}}}1
           {Ã}{{\~A}}1 {Á}{{\'A}}1 {É}{{\'E}}1 {Ê}{{\^E}}1 {Í}{{\'I}}1
           {Ó}{{\'O}}1 {Ô}{{\^O}}1 {Õ}{{\~O}}1 {Ú}{{\'U}}1 {Ç}{{\c{c}}}1
}
\lstdefinestyle{terminal}{
  basicstyle=\ttfamily\footnotesize\color{terminalfg},
  backgroundcolor=\color{terminalbg},
  frame=single,
  framesep=6pt,
  rulecolor=\color{deitelblue!50},
  breaklines=true,
  showstringspaces=false,
  literate={ã}{{\~a}}1 {á}{{\'a}}1 {é}{{\'e}}1 {ê}{{\^e}}1 {í}{{\'i}}1
           {ó}{{\'o}}1 {ô}{{\^o}}1 {õ}{{\~o}}1 {ú}{{\'u}}1 {ç}{{\c{c}}}1
}

\hypersetup{
  colorlinks=true,
  linkcolor=deitelblue,
  urlcolor=deitelblue,
  citecolor=deitelblue,
  pdftitle={C: Como Programar - Cap. 15 - Exercicios},
  pdfauthor={Resolucao didatica no estilo Deitel}
}

\pagestyle{fancy}
\fancyhf{}
\fancyhead[L]{\small\textit{C: Como Programar} --- Cap.~15}
\fancyhead[R]{\small Exerc\'icios}
\fancyfoot[C]{\small\thepage}
\renewcommand{\headrulewidth}{0.4pt}

\titleformat{\section}{\Large\bfseries\color{deitelblue}}{\thesection}{1em}{}
\titleformat{\subsection}{\large\bfseries\color{deitelblue}}{\thesubsection}{1em}{}

\newtcolorbox{boaprat}{
  enhanced, breakable, colback=boapratcolor!8, colframe=boapratcolor,
  fonttitle=\bfseries, title={\textbf{Boa Pr\'atica de Programa\c{c}\~ao}},
  left=8pt, right=8pt, top=4pt, bottom=4pt
}
\newtcolorbox{errocomum}{
  enhanced, breakable, colback=errocolor!8, colframe=errocolor,
  fonttitle=\bfseries, title={\textbf{Erro Comum de Programa\c{c}\~ao}},
  left=8pt, right=8pt, top=4pt, bottom=4pt
}
\newtcolorbox{obseng}{
  enhanced, breakable, colback=engcolor!8, colframe=engcolor,
  fonttitle=\bfseries, title={\textbf{Observa\c{c}\~ao de Engenharia de Software}},
  left=8pt, right=8pt, top=4pt, bottom=4pt
}
\newtcolorbox{dicadesemp}{
  enhanced, breakable, colback=desempcolor!8, colframe=desempcolor,
  fonttitle=\bfseries, title={\textbf{Dica de Desempenho}},
  left=8pt, right=8pt, top=4pt, bottom=4pt
}
\newtcolorbox{dicaportab}{
  enhanced, breakable, colback=portabcolor!8, colframe=portabcolor,
  fonttitle=\bfseries, title={\textbf{Dica de Portabilidade}},
  left=8pt, right=8pt, top=4pt, bottom=4pt
}
\newtcolorbox{enunciado}{
  enhanced, breakable, colback=deitelblue!5, colframe=deitelblue!70,
  boxrule=0.6pt, left=8pt, right=8pt, top=4pt, bottom=4pt
}
\newtcolorbox{saidabox}{
  enhanced, breakable, colback=gray!8, colframe=gray!50,
  boxrule=0.6pt, fonttitle=\bfseries\small,
  title={\textbf{Sa\'ida da execu\c{c}\~ao}},
  left=6pt, right=6pt, top=3pt, bottom=3pt
}

\title{
  \vspace{-1cm}
  {\Huge\bfseries\color{deitelblue} Cap\'itulo 15}\\[0.3em]
  {\Large\bfseries C++: um C melhor --- Introdu\c{c}\~ao \`a Tecnologia de Objeto}\\[0.8em]
  {\large\itshape Exerc\'icios --- Resolu\c{c}\~ao Did\'atica com C\'odigo C++}
}
\author{}
\date{}

\begin{document}
\maketitle
\thispagestyle{empty}
\vspace{-2em}

\begin{center}
\rule{0.85\textwidth}{0.5pt}\\[0.4em]
\textit{``Neste documento resolvemos os Exerc\'icios 15.5 a 15.10 do Cap\'itulo 15. Cada solu\c{c}\~ao apresenta o programa em C++ completo, comentado, seguido de an\'alise da arquitetura adotada, exemplos de execu\c{c}\~ao real e observa\c{c}\~oes did\'aticas. Todos os c\'odigos foram compilados com \texttt{g++ -Wall -Wextra -std=c++14} e executados sem warnings. Esta \'e a primeira incurs\~ao do leitor em c\'odigo C++ propriamente dito --- ainda na vertente \emph{procedural} da linguagem (refer\^encias, sobrecarga, argumentos default, templates, \texttt{inline}, escopo); as classes vir\~ao a partir do Cap\'itulo 16.''}\\[0.4em]
\rule{0.85\textwidth}{0.5pt}
\end{center}

\vspace{1em}

% ============================================================
\section{Exerc\'icio 15.5 --- \'Area de um C\'irculo com Fun\c{c}\~ao Inline}
% ============================================================

\begin{enunciado}
\textbf{Enunciado.} Escreva um programa em C++ que pe\c{c}a o raio de um c\'irculo e chame a fun\c{c}\~ao \texttt{inline} \texttt{circleArea} para calcular sua \'area.
\end{enunciado}

\subsection*{An\'alise do problema}

\'E o irm\~ao bidimensional do Exerc\'icio 15.4. A f\'ormula da \'area do c\'irculo \'e $A = \pi r^{2}$. Como antes, declaramos \texttt{circleArea} como \texttt{inline}: como o corpo \'e uma f\'ormula de uma linha, a expans\~ao em linha (\emph{inlining}) tem alto potencial de elimina\c{c}\~ao do custo de chamada de fun\c{c}\~ao.

\subsection*{C\'odigo em C++}

\lstinputlisting[style=cppcode]{codigos/ex15_05.cpp}

\subsection*{Execu\c{c}\~ao}

\begin{saidabox}
\begin{lstlisting}[style=terminal]
$ g++ -Wall -Wextra -std=c++14 ex15_05.cpp -o ex15_05
$ ./ex15_05
Digite o raio do circulo: 3
A area do circulo com raio 3 e 28.2743
\end{lstlisting}
\end{saidabox}

\subsection*{Discuss\~ao}

Conferindo: $A = \pi \cdot 3^{2} = 9\pi \approx 28{,}27$. Bate.

\begin{itemize}[leftmargin=*,itemsep=0pt]
  \item O par\^ametro \texttt{const double radius} indica que \texttt{circleArea} \emph{n\~ao} modificar\'a o raio recebido. Adicionar \texttt{const} a par\^ametros que voc\^e n\~ao pretende alterar \'e uma forma de documenta\c{c}\~ao executavel: o compilador valida que voc\^e n\~ao quebrou essa promessa.
  \item Como o par\^ametro \'e passado \emph{por valor}, o \texttt{const} aqui \'e meramente uma autodisciplina interna \`a fun\c{c}\~ao --- o chamador n\~ao \'e afetado.
  \item Outro ponto: \texttt{pow(radius, 2)} foi usado por simetria com 15.4, mas \texttt{radius * radius} seria mais r\'apido (uma multiplica\c{c}\~ao versus uma chamada de fun\c{c}\~ao gen\'erica).
\end{itemize}

\begin{dicadesemp}
\texttt{pow(x, 2)} \'e um c\'alculo gen\'erico que aceita expoentes arbitr\'arios em ponto flutuante --- internamente usa \texttt{exp(2 * log(x))} ou um caso especial. Para pot\^encias inteiras pequenas e conhecidas em tempo de compila\c{c}\~ao, prefira a multiplica\c{c}\~ao expl\'icita. Em c\'odigo de alto desempenho, essa diferen\c{c}a pode somar \emph{ordens de magnitude}.
\end{dicadesemp}

% ============================================================
\section{Exerc\'icio 15.6 --- \texttt{tripleByValue} vs \texttt{tripleByReference}}
% ============================================================

\begin{enunciado}
\textbf{Enunciado.} Escreva um programa em C++ com duas fun\c{c}\~oes alternativas, cada uma simplesmente triplicando a vari\'avel \texttt{count} definida em \texttt{main}, e compare-as:
\begin{enumerate}[label=(\alph*)]
  \item \texttt{tripleByValue}, que passa uma c\'opia de \texttt{count} por valor, triplica a c\'opia e retorna o novo valor;
  \item \texttt{tripleByReference}, que passa \texttt{count} por refer\^encia e triplica o valor original atrav\'es de seu \emph{alias}.
\end{enumerate}
\end{enunciado}

\subsection*{An\'alise do problema}

Este \'e o exerc\'icio-chave para fixar a diferen\c{c}a entre os dois mecanismos. A pergunta n\~ao \'e ``qual triplica?'' (ambos triplicam o valor 10 e produzem 30); \'e ``qual modifica o original?''. Vamos verificar empiricamente, imprimindo o valor de \texttt{count} \emph{antes} e \emph{depois} de cada chamada.

\subsection*{C\'odigo em C++}

\lstinputlisting[style=cppcode]{codigos/ex15_06.cpp}

\subsection*{Execu\c{c}\~ao}

\begin{saidabox}
\begin{lstlisting}[style=terminal]
$ ./ex15_06
Antes de chamar tripleByValue:     count = 10
Apos chamar tripleByValue:         count = 10  (retorno = 30)
  --> Observe que count NAO mudou.

Antes de chamar tripleByReference: count = 10
Apos chamar tripleByReference:     count = 30
  --> Observe que count FOI modificado.
\end{lstlisting}
\end{saidabox}

\subsection*{Discuss\~ao --- Compara\c{c}\~ao das duas t\'ecnicas}

\begin{center}
\renewcommand{\arraystretch}{1.4}
\begin{tabular}{|p{4cm}|p{4.8cm}|p{4.8cm}|}
\hline
\textbf{Aspecto} & \textbf{Por valor (\texttt{tripleByValue})} & \textbf{Por refer\^encia (\texttt{tripleByReference})} \\
\hline
Declara\c{c}\~ao         & \texttt{int tripleByValue(int x)} & \texttt{void tripleByReference(int \&x)} \\
Chamada                  & \texttt{r = tripleByValue(count);} & \texttt{tripleByReference(count);} \\
A vari\'avel original muda? & N\~ao, \texttt{x} \'e c\'opia & Sim, \texttt{x} \'e \emph{alias} \\
Necessita retornar valor? & Sim (sen\~ao perde-se o resultado) & N\~ao (modifica direto) \\
Seguran\c{c}a            & Melhor (n\~ao h\'a \emph{aliasing}) & Exige cuidado \\
Custo                    & Custo de c\'opia (irrelevante p/ \texttt{int}) & Sem c\'opia \\
\hline
\end{tabular}
\end{center}

\begin{obseng}
A passagem por valor \emph{preserva} o argumento original, o que \'e desej\'avel sempre que poss\'ivel: o c\'odigo torna-se mais f\'acil de raciocinar, porque a chamada n\~ao tem \emph{efeitos colaterais} sobre as vari\'aveis do chamador. A passagem por refer\^encia ``ganha tempo'' (sem c\'opia) e ``ganha expressividade'' (n\~ao precisa retornar via \texttt{return}), mas em troca o leitor do c\'odigo precisa lembrar que, ap\'os a chamada, a vari\'avel pode ter mudado.
\end{obseng}

\begin{boaprat}
Quando voc\^e \emph{l\^e} c\'odigo C++, o operador \texttt{\&} no prot\'otipo do par\^ametro \'e um \emph{sinal vis\'ivel} de que aquele argumento pode ser modificado. Costume-se a procurar esse sinal --- ele responde rapidamente \`a pergunta ``esta chamada vai mexer em alguma das minhas vari\'aveis?''.
\end{boaprat}

% ============================================================
\section{Exerc\'icio 15.7 --- Finalidade do Operador Un\'ario \texttt{::}}
% ============================================================

\begin{enunciado}
\textbf{Enunciado.} Qual \'e a finalidade do operador un\'ario de resolu\c{c}\~ao de escopo?
\end{enunciado}

\subsection*{Resposta}

\begin{enunciado}
O \textbf{operador un\'ario de resolu\c{c}\~ao de escopo} (\texttt{::}) permite acessar uma vari\'avel \emph{global} mesmo quando ela est\'a sendo ``ocultada'' por uma vari\'avel \emph{local} de mesmo nome no escopo atual. Sem o operador, qualquer refer\^encia ao identificador resolveria para a vari\'avel local (a mais pr\'oxima); com o prefixo \texttt{::}, o compilador faz a busca a partir do escopo \emph{global}.
\end{enunciado}

\subsection*{Demonstra\c{c}\~ao}

\begin{lstlisting}[style=cppcode]
#include <iostream>
using namespace std;

int valor = 100;        // global

int main() {
    int valor = 7;      // local; oculta a global
    cout << "Local:  " << valor   << endl;   // imprime 7
    cout << "Global: " << ::valor << endl;   // imprime 100
    return 0;
}
\end{lstlisting}

\subsection*{Discuss\~ao adicional}

H\'a ainda uma \emph{segunda} forma do operador \texttt{::}, n\~ao coberta por esta pergunta espec\'ifica mas onipresente em C++: o operador \textbf{bin\'ario} de resolu\c{c}\~ao de escopo, usado para qualificar nomes pertencentes a \emph{classes} ou \emph{\emph{namespaces}}. Por exemplo, \texttt{std::cout} significa ``o identificador \texttt{cout} dentro do \emph{namespace} \texttt{std}''. Ambas as formas usam o mesmo simbolo \texttt{::}, mas conceitualmente s\~ao distintas: a un\'aria sobe at\'e o escopo global; a bin\'aria desce at\'e um escopo nomeado.

\begin{obseng}
A exist\^encia do operador \texttt{::} un\'ario \'e um \emph{sintoma} de um problema, n\~ao a \emph{solu\c{c}\~ao} ideal. Boas pr\'aticas de engenharia de \emph{software} recomendam evitar nomes globais (use \emph{namespaces}, classes, ou par\^ametros) e, quando os globais s\~ao inevit\'aveis, dar-lhes nomes distintos dos locais. Em c\'odigo legado ou bibliotecas n\~ao-modulares, contudo, \texttt{::} ainda salva vidas.
\end{obseng}

% ============================================================
\section{Exerc\'icio 15.8 --- Template de Fun\c{c}\~ao \texttt{min}}
% ============================================================

\begin{enunciado}
\textbf{Enunciado.} Escreva um programa que use um template de fun\c{c}\~ao chamado \texttt{min} para determinar o menor de dois argumentos. Teste com inteiros, caracteres e ponto flutuante.
\end{enunciado}

\subsection*{An\'alise do problema}

A id\'eia da programa\c{c}\~ao gen\'erica: escrever \emph{uma s\'o} fun\c{c}\~ao que sirva para todos os tipos comparaveis. Para que isso seja poss\'ivel, o corpo da fun\c{c}\~ao s\'o pode usar opera\c{c}\~oes que existam para todos os tipos esperados. No nosso caso, basta o operador \texttt{<}, que est\'a definido para inteiros, ponto flutuante, caracteres e mesmo strings.

\begin{obseng}
Renomeamos a fun\c{c}\~ao para \texttt{minimo} no c\'odigo, evitando colis\~ao com a fun\c{c}\~ao \texttt{std::min} da biblioteca padr\~ao (que est\'a sempre dispon\'ivel quando se faz \texttt{using namespace std}). \'E uma quest\~ao puramente did\'atica --- em c\'odigo de produ\c{c}\~ao, simplesmente usar\'iamos \texttt{std::min}.
\end{obseng}

\subsection*{C\'odigo em C++}

\lstinputlisting[style=cppcode]{codigos/ex15_08.cpp}

\subsection*{Execu\c{c}\~ao}

\begin{saidabox}
\begin{lstlisting}[style=terminal]
$ ./ex15_08
min(47, 23) = 23
min(3.14, 2.71) = 2.71
min('M', 'A') = 'A'
\end{lstlisting}
\end{saidabox}

\subsection*{Discuss\~ao}

\begin{itemize}[leftmargin=*,itemsep=0pt]
  \item Para cada chamada, o compilador examina os tipos dos argumentos e \emph{instancia} uma vers\~ao espec\'ifica do template. Internamente, gera-se algo equivalente a tr\^es fun\c{c}\~oes: \texttt{minimo<int>}, \texttt{minimo<double>}, \texttt{minimo<char>}.
  \item A compara\c{c}\~ao de caracteres usa a ordem da tabela ASCII: \texttt{'A'} (65) \'e menor que \texttt{'M'} (77), por isso o resultado \'e \texttt{'A'}.
  \item Os par\^ametros foram passados por valor. Se quis\'essemos evitar c\'opias para tipos grandes, escrever\'iamos \texttt{const T \&a, const T \&b}.
\end{itemize}

\begin{boaprat}
A pr\'opria biblioteca padr\~ao oferece \texttt{std::min(a, b)} em \texttt{<algorithm>}. Em c\'odigo real, prefira-a. O exerc\'icio nos pede a \emph{implementa\c{c}\~ao} did\'atica, n\~ao a substitui\c{c}\~ao da biblioteca padr\~ao.
\end{boaprat}

\begin{errocomum}
Aplicar templates a tipos que n\~ao suportam todas as opera\c{c}\~oes usadas no corpo gera \emph{mensagens de erro de compila\c{c}\~ao extensas e dif\'iceis de ler}. Por exemplo, \texttt{minimo(structA, structB)}, se \texttt{structA} n\~ao definir \texttt{operator<}, produzir\'a um erro com dezenas de linhas. C++20 oferece \emph{concepts} para tornar essas mensagens muito mais claras.
\end{errocomum}

% ============================================================
\section{Exerc\'icio 15.9 --- Template de Fun\c{c}\~ao \texttt{max}}
% ============================================================

\begin{enunciado}
\textbf{Enunciado.} Escreva um programa que use um template de fun\c{c}\~ao chamado \texttt{max} para determinar o maior de dois argumentos. Teste com inteiros, caracteres e ponto flutuante.
\end{enunciado}

\subsection*{An\'alise do problema}

A solu\c{c}\~ao \'e o ``espelho'' do exerc\'icio anterior: substitu\'imos o operador \texttt{<} por \texttt{>}.

\subsection*{C\'odigo em C++}

\lstinputlisting[style=cppcode]{codigos/ex15_09.cpp}

\subsection*{Execu\c{c}\~ao}

\begin{saidabox}
\begin{lstlisting}[style=terminal]
$ ./ex15_09
max(47, 23) = 47
max(3.14, 2.71) = 3.14
max('M', 'A') = 'M'
\end{lstlisting}
\end{saidabox}

\subsection*{Discuss\~ao}

A simetria com 15.8 ilustra como templates removem a duplica\c{c}\~ao de c\'odigo no eixo do \emph{tipo}, mas ainda exigem c\'odigo separado quando muda a \emph{l\'ogica}. Se quis\'essemos parametrizar tamb\'em a opera\c{c}\~ao de compara\c{c}\~ao, ter\'iamos que passar uma fun\c{c}\~ao como argumento --- conceito visto em C com \emph{ponteiros para fun\c{c}\~ao} (Cap.~7), e que em C++ ganha forma elegante com \emph{functors} e express\~oes \emph{lambda} (Caps.~20 e seguintes).

\begin{obseng}
A pr\'opria biblioteca padr\~ao implementa \texttt{std::max} e \texttt{std::min} usando templates muito parecidos com os nossos --- s\'o que com sobrecargas adicionais para iterators, listas de inicializa\c{c}\~ao e fun\c{c}\~oes de compara\c{c}\~ao personalizadas. Isso \'e uma observa\c{c}\~ao consoladora: o que voc\^e escreveu aqui n\~ao \'e um esbo\c{c}o imaturo; \'e essencialmente o \emph{n\'ucleo} da implementa\c{c}\~ao real.
\end{obseng}

\newpage
% ============================================================
\section{Exerc\'icio 15.10 --- Caça aos Erros}
% ============================================================

\begin{enunciado}
\textbf{Enunciado.} Determine se os segmentos de programa a seguir cont\^em erros. Explique como cada erro pode ser corrigido. Para um segmento em particular, \'e poss\'ivel que nenhum erro esteja presente.
\end{enunciado}

% ----- (a) -----
\subsection*{(a) Template cujo par\^ametro de tipo n\~ao \'e usado}

\textbf{Segmento original:}
\begin{lstlisting}[style=cppcode,numbers=none]
template < class A >
int sum( int num1, int num2, int num3 )
{
    return num1 + num2 + num3;
}
\end{lstlisting}

\textbf{Diagn\'ostico.} O segmento \emph{compila}, mas cont\'em uma incoer\^encia conceitual: o par\^ametro de tipo \texttt{A} \'e declarado, por\'em \emph{nunca usado} no corpo da fun\c{c}\~ao --- todos os par\^ametros e o retorno s\~ao \texttt{int} fixos. Isso anula completamente a vantagem do template: a fun\c{c}\~ao \emph{n\~ao} \'e gen\'erica; ela aceita apenas tr\^es \texttt{int}s. Para invocar essa fun\c{c}\~ao, o programador ainda precisaria escrever \texttt{sum<algumTipo>(1,2,3)}, o que \'e in\'util.

\textbf{Corre\c{c}\~ao.} H\'a duas op\c{c}\~oes:

\textit{Op\c{c}\~ao 1 --- Tornar a fun\c{c}\~ao genuinamente gen\'erica:}
\begin{lstlisting}[style=cppcode,numbers=none]
template <typename A>
A sum(A num1, A num2, A num3) {
    return num1 + num2 + num3;
}
\end{lstlisting}

\textit{Op\c{c}\~ao 2 --- Remover o template:}
\begin{lstlisting}[style=cppcode,numbers=none]
int sum(int num1, int num2, int num3) {
    return num1 + num2 + num3;
}
\end{lstlisting}

A op\c{c}\~ao 1 \'e prefer\'ivel se a inten\c{c}\~ao original era de fato uma soma gen\'erica.

% ----- (b) -----
\subsection*{(b) Fun\c{c}\~ao \texttt{void} que retorna valor}

\textbf{Segmento original:}
\begin{lstlisting}[style=cppcode,numbers=none]
void printResults( int x, int y )
{
    cout << "A soma e " << x + y << '\n';
    return x + y;
}
\end{lstlisting}

\textbf{Diagn\'ostico.} Erro de compila\c{c}\~ao. A fun\c{c}\~ao foi declarada como \texttt{void} (n\~ao retorna nada), mas o corpo cont\'em \texttt{return x + y;} --- tenta-se retornar um \texttt{int}. O compilador rejeita essa contradi\c{c}\~ao.

\textbf{Corre\c{c}\~ao.} Decida o que a fun\c{c}\~ao deve fazer:

\textit{Op\c{c}\~ao 1 --- Se for para retornar a soma, mude o tipo de retorno:}
\begin{lstlisting}[style=cppcode,numbers=none]
int printResults(int x, int y) {
    cout << "A soma e " << x + y << '\n';
    return x + y;
}
\end{lstlisting}

\textit{Op\c{c}\~ao 2 --- Se for s\'o para imprimir, remova o \texttt{return}:}
\begin{lstlisting}[style=cppcode,numbers=none]
void printResults(int x, int y) {
    cout << "A soma e " << x + y << '\n';
}
\end{lstlisting}

% ----- (c) -----
\subsection*{(c) Template sem \texttt{class} ou \texttt{typename}}

\textbf{Segmento original:}
\begin{lstlisting}[style=cppcode,numbers=none]
template < A >
A product( A num1, A num2, A num3 )
{
    return num1 * num2 * num3;
}
\end{lstlisting}

\textbf{Diagn\'ostico.} Erro de sintaxe. Na declara\c{c}\~ao de um par\^ametro de tipo de template, deve-se usar a palavra-chave \texttt{class} ou \texttt{typename} \emph{antes} do identificador. Sem ela, o compilador interpreta \texttt{A} como o nome de uma vari\'avel de algum tipo (e n\~ao acha esse tipo) e protesta.

\textbf{Corre\c{c}\~ao.} Adicione \texttt{class} ou \texttt{typename}:

\begin{lstlisting}[style=cppcode,numbers=none]
template <class A>
A product(A num1, A num2, A num3) {
    return num1 * num2 * num3;
}
\end{lstlisting}

\noindent ou, equivalentemente, com \texttt{typename}:

\begin{lstlisting}[style=cppcode,numbers=none]
template <typename A>
A product(A num1, A num2, A num3) {
    return num1 * num2 * num3;
}
\end{lstlisting}

\begin{obseng}
\texttt{class} e \texttt{typename}, no contexto de par\^ametros de template, s\~ao \emph{intercambiaveis}. Hist\'oricamente, o C++ original usava apenas \texttt{class} (o que era confuso porque o par\^ametro n\~ao precisa ser uma classe --- pode ser qualquer tipo, inclusive primitivos). \texttt{typename} foi adicionado posteriormente para tornar a inten\c{c}\~ao clara, e \'e hoje a forma preferida. Mas \texttt{class} continua v\'alida por raz\~oes de compatibilidade.
\end{obseng}

% ----- (d) -----
\subsection*{(d) Sobrecarga apenas pelo tipo de retorno}

\textbf{Segmento original:}
\begin{lstlisting}[style=cppcode,numbers=none]
double cube( int );
int    cube( int );
\end{lstlisting}

\textbf{Diagn\'ostico.} Erro de compila\c{c}\~ao. Estas duas declara\c{c}\~oes tentam \emph{sobrecarregar} \texttt{cube}, mas as duas vers\~oes diferem \emph{apenas no tipo de retorno}, mantendo id\^enticas as listas de par\^ametros (\texttt{int}). A regra de C++ \'e clara: a sobrecarga distingue fun\c{c}\~oes pelas \emph{assinaturas}, e a assinatura \'e composta pelo \emph{nome} e pelos \emph{tipos dos par\^ametros} --- \emph{n\~ao} pelo tipo de retorno. Logo, do ponto de vista do compilador, essas duas fun\c{c}\~oes s\~ao \emph{a mesma}, e ele rejeita a duplicidade.

\textbf{Por que essa regra?} Imagine que o compilador encontre \texttt{int x = cube(5);}. Qual das duas escolher? A pelo tipo de retorno depende do contexto da \emph{chamada}, e isso introduz ambiguidade insol\'uvel em muitos cen\'arios. Para evitar, o C++ proibiu de sa\'ida.

\textbf{Corre\c{c}\~ao.} Algumas op\c{c}\~oes:

\textit{Op\c{c}\~ao 1 --- D\^e nomes distintos:}
\begin{lstlisting}[style=cppcode,numbers=none]
double cubeD(int);   // versao em ponto flutuante
int    cubeI(int);   // versao inteira
\end{lstlisting}

\textit{Op\c{c}\~ao 2 --- Diferencie pelos tipos de par\^ametros:}
\begin{lstlisting}[style=cppcode,numbers=none]
double cube(double);   // recebe double, retorna double
int    cube(int);      // recebe int,    retorna int
\end{lstlisting}

\textit{Op\c{c}\~ao 3 --- Use um template:}
\begin{lstlisting}[style=cppcode,numbers=none]
template <typename T>
T cube(T x) { return x * x * x; }
\end{lstlisting}

\subsection*{Programa de verifica\c{c}\~ao}

Compilamos um programa contendo as quatro vers\~oes corrigidas para confirmar que todas funcionam:

\lstinputlisting[style=cppcode]{codigos/ex15_10.cpp}

\subsection*{Execu\c{c}\~ao}

\begin{saidabox}
\begin{lstlisting}[style=terminal]
$ g++ -Wall -Wextra -std=c++14 ex15_10.cpp -o ex15_10
$ ./ex15_10
sum<int>(1,2,3)         = 6
sum<double>(1.5,2.5,3.0)= 7
A soma e 30
printResults retornou: 30
product<int>(2,3,4)     = 24
product<double>(1.5,2,4)= 12
cubeD(3) = 27  cubeI(4) = 64
\end{lstlisting}
\end{saidabox}

\subsection*{Resumo --- Tabela de diagn\'osticos}

\begin{center}
\renewcommand{\arraystretch}{1.4}
\begin{tabular}{|c|p{5cm}|p{6.5cm}|}
\hline
\textbf{Item} & \textbf{Erro} & \textbf{Corre\c{c}\~ao} \\
\hline
(a) & Par\^ametro de template n\~ao usado & Substituir \texttt{int} por \texttt{A}, ou remover \texttt{template} \\
(b) & \texttt{void} com \texttt{return x+y} & Mudar retorno para \texttt{int}, ou remover \texttt{return} \\
(c) & Falta \texttt{class}/\texttt{typename} & Inserir \texttt{class} ou \texttt{typename} antes de \texttt{A} \\
(d) & Sobrecarga apenas pelo retorno & Diferenciar nomes ou tipos de par\^ametros \\
\hline
\end{tabular}
\end{center}

\begin{obseng}
Caça-aos-erros como o (d) acima ensinam um princ\'ipio importante do projeto de linguagens: \emph{ambiguidade \'e o inimigo}. A escolha do C++ de \emph{n\~ao} permitir sobrecarga apenas pelo tipo de retorno preserva a regra de que toda chamada de fun\c{c}\~ao tem uma e somente uma resolu\c{c}\~ao --- regra que torna possi\'vel ao compilador (e ao leitor humano) entender o c\'odigo. Se esta regra fosse violada, todo \texttt{cube(5)} teria que ser examinado em seu contexto para decidir qual chamar.
\end{obseng}

% ============================================================
\section*{Encerramento}
\addcontentsline{toc}{section}{Encerramento}
% ============================================================

Os Exerc\'icios 15.5 a 15.10 oferecem o leitor um \emph{primeiro contato pr\'atico} com C++: usamos \texttt{cin}/\texttt{cout} no lugar de \texttt{scanf}/\texttt{printf}, conhecemos \texttt{inline}, refer\^encias e templates, e identificamos as armadilhas mais comuns na transi\c{c}\~ao a partir de C. Em todos os c\'odigos, contudo, permanecemos no paradigma \emph{procedural} --- fun\c{c}\~oes operando sobre dados.

A partir do Cap\'itulo 16 daremos o passo conceitual fundamental: introduziremos a \textbf{classe} como mecanismo de encapsulamento de \emph{dados} e \emph{comportamento}. As fun\c{c}\~oes deixar\~ao de ser entidades soltas e tornar-se-\~ao \emph{m\'etodos} pertencentes a um tipo. O leitor que dominou tudo at\'e aqui est\'a perfeitamente preparado para isso.

\vspace{1em}
\begin{center}\rule{0.5\textwidth}{0.4pt}\end{center}

\end{document}
